把两个模型 checkpoint 的权重相减、取范数，得到一个数，然后说这两个模型``差了这么多''。这个动作在工程里随处可见，可它量的是参数数值的差，而我们关心的一向是输出分布的差。同样一步 \(0.01\)，落在参数空间的不同位置上，对分布的影响可以相差几个数量级。本单元要做的事，是给参数空间换一把尺子。

换尺子的材料早就备齐了。Fisher 信息把``参数微小改变时分布有多容易区分''压成一个矩阵，黎曼度量把``长度''写进弯曲空间。信息几何把两者接起来：参数族本身是一个流形，Fisher 信息是它上面的度量。于是``两个分布有多不同''成了几何问题，而``该往哪个方向更新''成了这个几何里的最速下降问题。

\subsection{一个能立刻算出来的例子}

\(\Bernoulli(p)\) 族的 Fisher 信息是 \(I(p)=1/\bigl(p(1-p)\bigr)\)。在 \(p=0.5\) 处它等于 4，在 \(p=0.01\) 处约等于 101：靠近边界时，同样大小的参数扰动引起的分布改变要大得多，因为``几乎不发生''与``完全不发生''之间的差别很显眼。

这个数字直接决定了更新步长。自然梯度给出 \(\Delta p\propto I(p)^{-1}\nabla L=p(1-p)\nabla L\)，在 \(p\) 逼近 0 或 1 时自动收缩，而普通梯度下降在同一处会因为 \(\nabla L\) 的发散而失控。同一个 \(p(1-p)\) 出现在两端，两处的极端行为互相抵消——这是本单元反复出现的模式。

\subsection{三篇按对象、方向、局部距离推进}

第一篇把参数族看成几何对象。每个参数值是流形上的一个点，Fisher 信息在每个点上定义一个内积，由此得到的距离衡量分布的可区分程度而非参数的数值差。关键性质是坐标不变：换一套参数化，矩阵的元素变了，算出来的长度没变。

第二篇换掉梯度。把信任域的形状从欧氏球改成 Fisher 椭圆，最速下降方向随之变成 \(I(\theta)^{-1}\nabla L\)，即自然梯度。它对参数化方式不变，与 Newton 法的仿射不变性同出一源，区别在于预条件矩阵来自模型族还是来自目标函数的曲率。实际使用的都是近似——对角近似（Adam 大致落在这一类）、Kronecker 分解（K-FAC）——近似有多粗，不变性就损失多少。

第三篇把 KL 散度接进来。KL 在参数邻域内的二阶展开，其二次型系数正是 Fisher 矩阵，一阶项因 score 期望为零而消失。这一条把信息论、统计与优化连成一条线，也解释了强化学习信任域方法里那个 KL 约束的几何含义：它约束的是一个随位置改变形状的椭圆。

\begin{itemize}
\tightlist
\item \hyperref[sec:12-Real-Analysis-Support/09-Information-Geometry/01]{``统计流形把参数族变成几何对象''}；
\item \hyperref[sec:12-Real-Analysis-Support/09-Information-Geometry/02]{``自然梯度沿分布几何下降''}；
\item \hyperref[sec:12-Real-Analysis-Support/09-Information-Geometry/03]{``KL 散度与 Fisher 度量的局部关系''}。
\end{itemize}

顺序读下来最省力，因为后两篇都在用第一篇的度量。若手边的困惑具体到``为什么换个参数化训练曲线就变了''，可以直接从第二篇进入，再回头补几何的定义。
